\documentclass{zhvt-classic}

% 使用本地源流明体（GenRyuMin2 TW）按文件路径加载，不依赖系统字体名
% 版框（内框／界栏／书口侧栏／鱼尾）改用与版心印章同一印泥棕色：
%   印章图片 yshs.png 实色 = RGB(135,68,52) = HTML 874434
\usepackage{xcolor}
\usepackage{tikz}
% sealbrown 双注册：xcolor 供 TikZ/pgf，L3 供 cls 版框
\definecolor{sealbrown}{HTML}{874434}
\ExplSyntaxOn
\color_set:nnn { sealbrown } { HTML } { 874434 }
\ExplSyntaxOff

\zhvtset{
  grid color       = sealbrown,
  main font        = GenRyuMin2TW-SB,
  main font path   = C:/Users/Administrator/texmf/fonts/opentype/public/GenRyuMin2TW/,
jiazhu font      = GenRyuMin2TW-M,
  jiazhu font path = C:/Users/Administrator/texmf/fonts/opentype/public/GenRyuMin2TW/,
  % 页码框容纳 4 字：正常页「九十七」+前置空 1 格=4 行恰好；
  % 100+ 页「一百零一」共 5 行超出，末字「一/二/三」被框裁掉（2026-10-04 实测），故放宽到 5
  shukou page number chars = 5,
  % 书口书名框高从默认 6 字减到 4 字，把省出的 2 字高拨给 chapter 框，
  % 使 6 字书名+卷名放进 9 字高的 chapter 框后不再压线（2026-10-05）
  shukou title chars = 4,
}

\title{史记}[齐悼惠王世家]
\author{汉·司马迁}
\maker{闲叟漫记·重製}


% ===== 索隐述赞样式（2026-11-08 四次定稿 · 愚大人亲手改定参数）=====
% ⚠ 竖排前提：自定义字体家族必须带 RawFeature={vertical:+vert}；
%   simhei.ttf 加 +vhal 会整体加载失败静默回退，只写 +vert；simfang 同理。
\newCJKfontfamily{\shuzanfont}[Path=C:/Windows/Fonts/,RawFeature={vertical:+vert}]{simhei.ttf}
\newCJKfontfamily{\shuzanbody}[Path=C:/Windows/Fonts/,RawFeature={vertical:+vert}]{simfang.ttf}
% 题头：152x30pt 圆角 6pt 棕底白字；baseline (0,2) 使整框落入界行格子不压版框线；
% 白字节点 x 不补偿、y +8.2pt（像素级墨迹实测 dx/dy≈0，L/R 留白约 20pt）。
\newcommand{\shuzanhead}[1]{%
  \begin{tikzpicture}[baseline={(0,2)}]
    \fill[sealbrown, rounded corners=6pt]
      (-76pt,-15pt) rectangle (76pt,15pt);
    \node[text=white, inner sep=0pt, anchor=center] at (0pt,8.2pt)
      {\shuzanfont\fontsize{23}{23}\selectfont #1};
  \end{tikzpicture}%
}
% ★ 修正书口卷名位置：chapter 框起始位置改回 cls 默认值，避免全局影响。
%   只在每个卷首页通过 \juan 命令里的空命令微调该页位置（安全，不会跨页污染）。
\ExplSyntaxOn
\AtBeginDocument{
  \dim_set:Nn \l__zhvt_shukou_chapter_title_start_dim
    { \l__zhvt_shukou_title_height_dim + \l__zhvt_font_size_dim * 2 }
  \dim_set:Nn \l__zhvt_shukou_page_number_start_dim
    { \l__zhvt_shukou_chapter_title_start_dim + \l__zhvt_shukou_chapter_title_height_dim }
}
\ExplSyntaxOff

\begin{document}

% ★★ 版心（书口）排布覆写：把书名放到鱼尾之下（2026-10-04）
%
%   cls 默认版心自上而下：书名（title 静态框，在鱼尾之上）→ 鱼尾 → 篇名（chapter 动态框）→ 页码
%   本模板改为：          鱼尾（位置不动）→ 书名 → 篇名 → 页码分割线 → 页码
%
%   手法：把 cls 放在 title 静态框里的书名，改放进 chapter 动态框（鱼尾之下）；
%         title 静态框留空。书名取 \l__zhvt_title_str，即 \title{...} 的第一参数，
%         所以换书时只改 \title，不必在此重复写书名。
%
%   ⚠ 三条必须遵守的写法（2026-10-04 实测）：
%     ① 汉字必须用 \hbox{...} 包住。flowfram 会把框内容装进一个宽 14.4mm（= 一个 line_sep
%        = 书口栏宽）的纵向 minipage 并按该宽度折行，而整页靠 PDF /Rotate 90 显示、
%        框内文字本身不变向 —— 不包 \hbox 时「围炉夜话」4 字会被逐字折成 4 行，
%        在版面上就成了横排 4 字并横向溢出版框。
%     ② 竖排下**一个 \hbox 就是一整列**。书名与篇名要同处一列，必须放进同一个 \hbox；
%        分成两个 \hbox 时，后一个会被当成新的一列、甩到隔壁栏去
%        （实测「序」跑到书口栏外侧 14.4mm 的正文列上）。中间的 \quad 在框内是横向空白，
%        显示出来正好是书名与篇名之间的纵向间隔。
%     ③ ★ 内容盒要写成 **与框等宽的定宽盒** `\hbox to \l__zhvt_shukou_chapter_title_height_dim{…\hfil}`。
%        实测：chapter 框是把内容盒沿纵向**居中**摆放的，内容宽度一变，书名就跟着漂：
%        自然宽度的 \hbox（书名 38.4mm）落在 x=157.2mm，加篇名后（57.6mm）落在 x=152.4mm，
%        目录页（无篇名）与正文页的书名就差半格。做成等宽盒后，无论有没有篇名，
%        内容固定从框顶起排，实测两页均落在 x=142.8mm，不再漂。
\ExplSyntaxOn
% 把书口 chapter 框（书名+篇名）起始位置下移 1 字高，避免 4 字卷名顶到鱼尾
% ⚠ cls 的 \__zhvt_calculate_dimensions: 在 \AtBeginDocument 里执行，会重算该 dim；
%   必须把覆写挂到 begin document 之后，否则会被二次覆盖（2026-10-05 实测）。
% (旧方案已移到导言区，直接修改 chapter 框起始位置)
% 卷标题命令 #1=书口篇名(鱼尾) #2=正文卷名（含卷序，卷名内空格用 ~）
%
% ★★ 目录页两项修订（2026-10-05，愚大人指出）★★
%   ① 页码起点统一对齐：
%      原状 —— 目录项写成 \par\hspace{4em}卷名\hspace{4em}页码，页码起点 = 卷名实宽 + 4em。
%      卷名实宽随字数变：「卷一~梁惠王上」6 字 +~、「卷七~离娄上」5 字 +~，
%      差 1 个字位 → 卷七~卷十 4 条的页码整体比其余 10 条高 27.21pt（一个字），看着上下不齐。
%      改法 —— 卷名补全角空格（宽 = 1 字 = 27.215pt）补足到 **7 字符位**再写进目录，
%      14 条页码起点遂全部落在同一条竖线上（起点对齐；末字随页码字数自然参差）。
%      ⚠ 只对目录里的副本 \l_tmpc_tl 补，正文卷名仍用 \l_tmpa_tl 原样 —— 正文像素不受影响。
%   ② 卷名整体上移半字位：前导 \hspace 由 4em 改 3.5em。
%      实测原状目录块在版框内偏下（天头留白 139.6pt / 地脚留白 110.7pt），
%      上移 0.5em = 13.6pt 后上下基本均分，视觉居中。
\NewDocumentCommand { \juan } { mm } {
	\clearpage
	\tl_set:Nn \l__zhvt_chapter_title_tl { #1 }
	\tl_set:Nn \l_tmpa_tl { #2 }
	\tl_set_eq:NN \l_tmpc_tl \l_tmpa_tl
	\int_step_inline:nn { 7 - \str_count:n { #2 } }
	  { \tl_put_right:Nn \l_tmpc_tl { 　 } }
	\tl_set:Nx \l_tmpb_tl { \zhnum { page } }
	\tl_put_left:Nn \l_tmpb_tl { \hspace { 4em } }
	\tl_put_left:Nn \l_tmpb_tl { \l_tmpc_tl }
	\tl_put_left:Nn \l_tmpb_tl { \par \hspace { 3.5em } }
	\addtocontents { toc } { \l_tmpb_tl }
	\skip_horizontal:n { \l__zhvt_font_size_dim * 4 }
	\l_tmpa_tl
	\par
	\thispagestyle { empty }
}
\cs_set_protected:Npn \__zhvt_set_shukou_content: {
	% ★ 鱼尾上方留白（title 静态框）嵌入闲叟钤印：
	%   整页 /Rotate 90，故此处把图用 \rotatebox{90} 转正为竖章；
	%   尺寸卡在版心条幅内（书内约 7mm 宽 × 33.8mm 高，留余量不压鱼尾）。
	\setstaticcontents * { title } {
		\vbox_to_ht:nn { \l__zhvt_baseline_skip_dim } {
			\vskip 14.2pt
			\hbox_to_wd:nn { \l__zhvt_shukou_title_height_dim } {
				\hfil
				\rotatebox [ origin = c ] { 90 } { \includegraphics [ width = 7mm , height = 33.8mm ] { yshs.png } }
				\hfil
			}
		}
	}
% ★ 版心（书口）统一改法（2026-11-16 愚大人定）：18pt＋框内居中＋书口文字同版框色
%   ① 原 27.2pt 时 7 字长篇名（如「廉颇蔺相如列传」）末字溢出 chapter 框、压住页码；
%      改 \fontsize{18pt}{18pt}，前后 \hfil 令字串在框内居中（不随篇名字数漂移）。
%   ② \vbox{\vskip 8.2pt} 定量下移，抵消缩字号后栏宽方向 8.2pt 上偏（像素实测 y 中心 543.3 = 原版）。
%      ⚠ \raisebox / \vspace* 在 flowfram 框内无效；换字号须同步换补偿：16pt→10.1、19pt→7.4。
%   ③ chapter 与 page~number 均置 \color{sealbrown}（HTML 874434），与版框／界行同色；正文仍黑。
%   ⚠ 页码字位、钤印、鱼尾、界行、正文一律未动；五部合订本同此口径。
	\setdynamiccontents * { chapter }
	{ \color{sealbrown} \vbox { \vskip 8.2pt \hbox to \l__zhvt_shukou_chapter_title_height_dim
		{ \hfil \fontsize{18pt}{18pt}\selectfont \l__zhvt_title_str \quad \l__zhvt_chapter_title_tl \hfil } } }
	\setdynamiccontents * { page~number } { \color{sealbrown} \quad \zhnum { page } }
}
\__zhvt_set_shukou_content:
\ExplSyntaxOff

\maketitle

% 单篇，无多卷目录，不渲染目录页

\juan{齐悼惠王世家}{卷五十二~齐悼惠王世家第二十二}
\hspace*{2em}齐悼惠王刘肥者。高祖长庶男也。其母外妇也。曰曹氏。高祖六年。立肥为齐王。食七十城。诸民能齐言者皆予齐王。\\
\hspace*{2em}齐王。孝惠帝兄也。孝惠帝二年。齐王入朝。惠帝与齐王燕饮。亢礼如家人。吕太后怒。且诛齐王。齐王惧不得脱。乃用其内史勋计。献城阳郡。以为鲁元公主汤沐邑。吕太后喜。乃得辞就国。\\
\hspace*{2em}悼惠王即位十三年。以惠帝六年卒。子襄立。是为哀王。\\
\hspace*{2em}哀王元年。孝惠帝崩。吕太后称制。天下事皆决於高后。二年。高后立其兄子郦侯吕台为吕王。割齐之济南郡为吕王奉邑。\\
\hspace*{2em}哀王三年。其弟章入宿卫於汉。吕太后封为硃虚侯。以吕禄女妻之。後四年。封章弟兴居为东牟侯。皆宿卫长安中。\\
\hspace*{2em}哀王八年。高后割齐琅邪郡立营陵侯刘泽为琅邪王。\\
\hspace*{2em}其明年。赵王友入朝。幽死于邸。三赵王皆废。高后立诸吕诸吕为三王。擅权用事。\\
\hspace*{2em}硃虚侯年二十。有气力。忿刘氏不得职。尝入待高后燕饮。高后令硃虚侯刘章为酒吏。章自请曰。﹁臣。将种也。请得以军法行酒。﹂高后曰。﹁可。﹂酒酣。章进饮歌舞。已而曰。﹁请为太后言耕田歌。﹂高后兒子畜之。笑曰。﹁顾而父知田耳。若生而为王子。安知田乎。﹂章曰。﹁臣知之。﹂太后曰。﹁试为我言田。﹂章曰。﹁深耕穊种。立苗欲疏。非其种者。鉏而去之。﹂吕后默然。顷之。诸吕有一人醉。亡酒。章追。拔剑斩之。而还报曰。﹁有亡酒一人。臣谨行法斩之。﹂太后左右皆大惊。业已许其军法。无以罪也。因罢。自是之後。诸吕惮硃虚侯。虽大臣皆依硃虚侯。刘氏为益彊。\\
\hspace*{2em}其明年。高后崩。赵王吕禄为上将军。吕王产为相国。皆居长安中。聚兵以威大臣。欲为乱。硃虚侯章以吕禄女为妇。知其谋。乃使人阴出告其兄齐王。欲令发兵西。硃虚侯。东牟侯为内应。以诛诸吕。因立齐王为帝。\\
\hspace*{2em}齐王既闻此计。乃与其舅父驷钧。郎中令祝午。中尉魏勃阴谋发兵。齐相召平闻之。乃发卒卫王宫。魏勃绐召平曰。﹁王欲发兵。非有汉虎符验也。而相君围王。固善。勃请为君将兵卫卫王。﹂召平信之。乃使魏勃将兵围王宫。勃既将兵。使围相府。召平曰。﹁嗟乎。道家之言﹃当断不断。反受其乱﹄。乃是也。﹂遂自杀。於是齐王以驷钧为相。魏勃为将军。祝午为内史。悉发国中兵。使祝午东诈琅邪王曰。﹁吕氏作乱。齐王发兵欲西诛之。齐王自以兒子。年少。不习兵革之事。原举国委大王。大王自高帝将也。习战事。齐王不敢离兵。使臣请大王幸之临菑见齐王计事。并将齐兵以西平关中之乱。﹂琅邪王信之。以为然。驰见齐王。齐王与魏勃等因留琅邪王。而使祝午尽发琅邪国而并将其兵。\\
\hspace*{2em}琅邪王刘泽既见欺。不得反国。乃说齐王曰。﹁齐悼惠王高皇帝长子。推本言之。而大王高皇帝適长孙也。当立。今诸大臣狐疑未有所定。而泽於刘氏最为长年。大臣固待泽决计。今大王留臣无为也。不如使我入关计事。﹂齐王以为然。乃益具车送琅邪王。\\
\hspace*{2em}琅邪王既行。齐遂举兵西攻吕国之济南。於是齐哀王遗诸侯王书曰。﹁高帝平定天下。王诸子弟。悼惠王於齐。悼惠王薨。惠帝使留侯张良立臣为齐王。惠帝崩。高后用事。春秋高。听诸吕擅废高帝所立。又杀三赵王。灭梁。燕。赵以王诸吕。分齐国为四。忠臣进谏。上惑乱不听。今高后崩。皇帝春秋富。未能治天下。固恃大臣诸侯。今诸吕又擅自尊官。聚兵严威。劫列侯忠臣。矫制以令天下。宗庙所以危。今寡人率兵入诛不当为王者。﹂\\
\hspace*{2em}汉闻齐发兵而西。相国吕产乃遣大将军灌婴东击之。灌婴至荥阳。乃谋曰。﹁诸吕将兵居关中。欲危刘氏而自立。我今破齐还报。是益吕氏资也。﹂乃留兵屯荥阳。使使喻齐王及诸侯。与连和。以待吕氏之变而共诛之。齐王闻之。乃西取其故济南郡。亦屯兵於齐西界以待约。\\
\hspace*{2em}吕禄。吕产欲作乱关中。硃虚侯与太尉勃。丞相平等诛之。硃虚侯首先斩吕产。於是太尉勃等乃得尽诛诸吕。而琅邪王亦从齐至长安。\\
\hspace*{2em}大臣议欲立齐王。而琅邪王及大臣曰。﹁齐王母家驷钧。恶戾。虎而冠者也。方以吕氏故几乱天下。今又立齐王。是欲复为吕氏也。代王母家薄氏。君子长者。且代王又亲高帝子。於今见在。且最为长。以子则顺。以善人则大臣安。﹂於是大臣乃谋迎立代王。而遣硃虚侯以诛吕氏事告齐王。令罢兵。\\
\hspace*{2em}灌婴在荥阳。闻魏勃本教齐王反。既诛吕氏。罢齐兵。使使召责问魏勃。勃曰。﹁失火之家。岂暇先言大人而後救火乎。﹂因退立。股战而栗。恐不能言者。终无他语。灌将军熟视笑曰。﹁人谓魏勃勇。妄庸人耳。何能为乎。﹂乃罢魏勃。魏勃父以善鼓琴见秦皇帝。及魏勃少时。欲求见齐相曹参。家贫无以自通。乃常独早夜埽齐相舍人门外。相舍人怪之。以为物。而伺之。得勃。勃曰。﹁原见相君。无因。故为子埽。欲以求见。﹂於是舍人见勃曹参。因以为舍人。一为参御。言事。参以为贤。言之齐悼惠王。悼惠王召见。则拜为内史。始。悼惠王得自置二千石。及悼惠王卒而哀王立。勃用事。重於齐相。\\
\hspace*{2em}王既罢兵归。而代王来立。是为孝文帝。\\
\hspace*{2em}孝文帝元年。尽以高后时所割齐之城阳。琅邪。济南郡复与齐。而徙琅邪王王燕。益封硃虚侯。东牟侯各二千户。\\
\hspace*{2em}是岁。齐哀王卒。太子立。是为文王。\\
\hspace*{2em}齐文王元年。汉以齐之城阳郡立硃虚侯为城阳王。以齐济北郡立东牟侯为济北王。\\
\hspace*{2em}二年。济北王反。汉诛杀之。地入于汉。\\
\hspace*{2em}後二年。孝文帝尽封齐悼惠王子罢军等七人皆为列侯。\\
\hspace*{2em}齐文王立十四年卒。无子。国除。地入于汉。\\
\hspace*{2em}後一岁。孝文帝以所封悼惠王子分齐为王。齐孝王将闾以悼惠王子杨虚侯为齐王。故齐别郡尽以王悼惠王子。子志为济北王。子辟光为济南王。子贤为菑川王。子卬为胶西王。子雄渠为胶东王。与城阳。齐凡七王。\\
\hspace*{2em}齐孝王十一年。吴王濞。楚王戊反。兴兵西。告诸侯曰﹁将诛汉贼臣晁错以安宗庙﹂。胶西。胶东。菑川。济南皆擅发兵应吴楚。欲与齐。齐孝王狐疑。城守不听。三国兵共围齐。齐王使路中大夫告於天子。天子复令路中大夫还告齐王。﹁善坚守。吾兵今破吴楚矣。﹂路中大夫至。三国兵围临菑数重。无从入。三国将劫与路中大夫盟。曰。﹁若反言汉已破矣。齐趣下三国。不且见屠。﹂路中大夫既许之。至城下。望见齐王。曰。﹁汉已发兵百万。使太尉周亚夫击破吴楚。方引兵救齐。齐必坚守无下。﹂三国将诛路中大夫。\\
\hspace*{2em}齐初围急。阴与三国通谋。约未定。会闻路中大夫从汉来。喜。及其大臣乃复劝王毋下三国。居无何。汉将栾布。平阳侯等兵至齐。击破三国兵。解齐围。已而复闻齐初与三国有谋。将欲移兵伐齐。齐孝王惧。乃饮药自杀。景帝闻之。以为齐首善。以迫劫有谋。非其罪也。乃立孝王太子寿为齐王。是为懿王。续齐後。而胶西。胶东。济南。菑川王咸诛灭。地入于汉。徙济北王王菑川。齐懿王立二十二年卒。子次景立。是为厉王。\\
\hspace*{2em}齐厉王。其母曰纪太后。太后取其弟纪氏女为厉王后。王不爱纪氏女。太后欲其家重宠。令其长女纪翁主入王宫。正其後宫。毋令得近王。欲令爱纪氏女。王因与其姊翁主奸。\\
\hspace*{2em}齐有宦者徐甲。入事汉皇太后。皇太后有爱女曰脩成君。脩成君非刘氏。太后怜之。脩成君有女名娥。太后欲嫁之於诸侯。宦者甲乃请使齐。必令王上书请娥。皇太后喜。使甲之齐。是时齐人主父偃知甲之使齐以取后事。亦因谓甲。﹁即事成。幸言偃女原得充王後宫。﹂甲既至齐。风以此事。纪太后大怒。曰。﹁王有后。後宫具备。且甲。齐贫人。急乃为宦者。入事汉。无补益。乃欲乱吾王家。且主父偃何为者。乃欲以女充後宫。﹂徐甲大穷。还报皇太后曰。﹁王已原尚娥。然有一害。恐如燕王。﹂燕王者。与其子昆弟奸。新坐以死。亡国。故以燕感太后。太后曰。﹁无复言嫁女齐事。﹂事浸浔闻於天子。主父偃由此亦与齐有卻。\\
\hspace*{2em}主父偃方幸於天子。用事。因言。﹁齐临菑十万户。市租千金。人众殷富。巨於长安。此非天子亲弟爱子不得王此。今齐王於亲属益疏。﹂乃从容言。﹁吕太后时齐欲反。吴楚时孝王几为乱。今闻齐王与其姊乱。﹂於是天子乃拜主父偃为齐相。且正其事。主父偃既至齐。乃急治王後宫宦者为王通於姊翁主所者。令其辞证皆引王。王年少。惧大罪为吏所执诛。乃饮药自杀。绝无後。\\
\hspace*{2em}是时赵王惧主父偃一出废齐。恐其渐疏骨肉。乃上书言偃受金及轻重之短。天子亦既囚偃。公孙弘言。﹁齐王以忧死毋後。国入汉。非诛偃无以塞天下之望。﹂遂诛偃。\\
\hspace*{2em}齐厉王立五年死。毋後。国入于汉。\\
\hspace*{2em}齐悼惠王後尚有二国。城阳及菑川。菑川地比齐。天子怜齐。为悼惠王冢园在郡。割临菑东环悼惠王冢园邑尽以予菑川。以奉悼惠王祭祀。\\
\hspace*{2em}城阳景王章。齐悼惠王子。以硃虚侯与大臣共诛诸吕。而章身首先斩相国吕王产於未央宫。孝文帝既立。益封章二千户。赐金千斤。孝文二年。以齐之城阳郡立章为城阳王。立二年卒。子喜立。是为共王。\\
\hspace*{2em}共王八年。徙王淮南。四年。复还王城阳。凡三十三年卒。子延立。是为顷王。\\
\hspace*{2em}顷王二十年卒。子义立。是为敬王。敬王九年卒。子武立。是为惠王。惠王十一年卒。子顺立。是为荒王。荒王四十六年卒。子恢立。是为戴王。戴王八年卒。子景立。至建始三年。十五岁。卒。\\
\hspace*{2em}济北王兴居。齐悼惠王子。以东牟侯助大臣诛诸吕。功少。及文帝从代来。兴居曰。﹁请与太仆婴入清宫。﹂废少帝。共与大臣尊立孝文帝。\\
\hspace*{2em}孝文帝二年。以齐之济北郡立兴居为济北王。与城阳王俱立。立二年。反。始大臣诛吕氏时。硃虚侯功尤大。许尽以赵地王硃虚侯。尽以梁地王东牟侯。及孝文帝立。闻硃虚。东牟之初欲立齐王。故绌其功。及二年。王诸子。乃割齐二郡以王章。兴居。章。兴居自以失职夺功。章死。而兴居闻匈奴大入汉。汉多发兵。使丞相灌婴击之。文帝亲幸太原。以为天子自击胡。遂发兵反於济北。天子闻之。罢丞相及行兵。皆归长安。使棘蒲侯柴将军击破虏济北王。王自杀。地入于汉。为郡。\\
\hspace*{2em}後十年。文帝十六年。复以齐悼惠王子安都侯志为济北王。十一年。吴楚反时。志坚守。不与诸侯合谋。吴楚已平。徙志王菑川。\\
\hspace*{2em}济南王辟光。齐悼惠王子。以勒侯孝文十六年为济南王。十一年。与吴楚反。汉击破。杀辟光。以济南为郡。地入于汉。\\
\hspace*{2em}菑川王贤。齐悼惠王子。以武城侯文帝十六年为菑川王。十一年。与吴楚反。汉击破。杀贤。\\
\hspace*{2em}天子因徙济北王志王菑川。志亦齐悼惠王子。以安都侯王济北。菑川王反。毋後。乃徙济北王王菑川。凡立三十五年卒。谥为懿王。子建代立。是为靖王。二十年卒。子遗代立。是为顷王。三十六年卒。子终古立。是为思王。二十八年卒。子尚立。是为孝王。五年卒。子横立。至建始三年。十一岁。卒。\\
\hspace*{2em}胶西王卬。齐悼惠王子。以昌平侯文帝十六年为胶西王。十一年。与吴楚反。汉击破。杀卬。地入于汉。为胶西郡。\\
\hspace*{2em}胶东王雄渠。齐悼惠王子。以白石侯文帝十六年为胶东王。十一年。与吴楚反。汉击破。杀雄渠。地入于汉。为胶东郡。\\
\hspace*{2em}太史公曰。诸侯大国无过齐悼惠王。以海内初定。子弟少。激秦之无尺土封。故大封同姓。以填万民之心。及後分裂。固其理也。\\

% ===== 索隐述赞（司马贞《史记索隐》）· 统一样式 =====
% 题头「索隐述赞曰」五字并入反白框；正文每四句一列，定宽盒 \hfil 居列中；
% \fontsize{24}{41} 列距对齐全书界行（121+41k）；\clearpage 使述赞独占一页。
\clearpage
\shuzanhead{索隐述赞曰}\linebreak
{\shuzanbody\fontsize{24}{41}\selectfont
\hbox to \hsize{\hfil 汉矫秦制。树屏自彊。表海大国。悉封齐王。\hfil}\linebreak
\hbox to \hsize{\hfil 吕后肆怒。乃献城阳。哀王嗣立。其力不量。\hfil}\linebreak
\hbox to \hsize{\hfil 硃虚仕汉。功大策长。东牟受赏。称乱贻殃。\hfil}\linebreak
\hbox to \hsize{\hfil 胶东。济北。雄渠。辟光。\hfil}\linebreak
\hbox to \hsize{\hfil 齐虽七国。忠孝者昌。\hfil}%
\par}
\end{document}
